\documentclass[a4paper, 12pt]{article}

\usepackage{utils}

\renewcommand*{\today}{12 March 2026}

\begin{document}

\hotbox{Géometrie 2}{CM 5}{\today}

\section{Les angles}

\subsection{My name is BOND, James BOND}

\begin{definition}[(Orientation d'un espace vectoriel)]
    Soit $E$ un $\R$-espace vectoriel de dimension finie.

    \begin{itemize}
        \item Deux bases $\mathcal{B}$ et $\mathcal{B}'$ de $E$ ont la \textbf{même orientation} si la matrice de passage de $\mathcal{B}$ à $\mathcal{B}'$ a un déterminant strictement positif.
        \item La relation "avoir la même orientation" est une relation d'équivalence sur l'ensemble des bases de $E$. Il y a exactement deux classes d'équivalence, appelées les \textbf{orientations} de $E$.
        \item Orienter $E$ consiste à choisir arbitrairement une orientation de $E$, c'est à dire une classe d'équivalence de bases de $E$.
        Les bases appartenant à l'orientation choisie sont dites \textbf{bases directes}, les autres sont dites \textbf{bases indirectes}.
    \end{itemize}
\end{definition}

\begin{definition}[(Base orthonormée directe ou BOND)]
    Soit $E$ un espace euclidien orienté de dimension finie.

    Une base $\mathcal{B}$ de $E$ est une \textbf{base orthonormée directe} (BOND) si elle verifie les conditions
    \begin{itemize}
        \item $\mathcal{B}$ est une base orthonormée de $E$.
        \item $\mathcal{B}$ est une base directe de $E$.
    \end{itemize}
\end{definition}

\subsection{Angle de vecteurs}

\begin{theoreme}{}{angle_vecteurs}
    Soit $\mathcal{P}$ un plan vectoriel euclidien orienté.

    Soient $u, v \in \mathcal{P}$ tels que $\norm{u} = \norm{v} = 1$.

    \vspace{1cm}
    En notant $(a, b)$ les coordonnées de $v$ dans la BOND de premier vecteur $u$,
    il existe un unique $\theta \in \R$ modulo $2\pi$ tel que $a = \cos(\theta)$ et $b = \sin(\theta)$.
\end{theoreme}

\begin{demonstration}
    Il existe un unique $u^\perp \in \mathcal{P}$ tel que $\mathcal{B} = (u, u^\perp)$ soit une BOND de $\mathcal{P}$.

    On considère l'unique rotation vectorielle $r$ telle que $r(u) = v$.

    Ainsi on peut écrire
    $$
    Mat_{\mathcal{B}}(r) = \begin{pmatrix}
        a & -b \\
        b & a
    \end{pmatrix}
    $$

    Une rotation vectorielle conserve l'orientation et les distances, donc $det(r) = 1$ et
    $$
    a^2 + b^2 = 1
    $$
    Il existe un unique $\theta \in \R$ modulo $2\pi$ tel que $a = \cos(\theta)$ et $b = \sin(\theta)$.
\end{demonstration}

\begin{definition}[(Angle de vecteurs)]
    Soit $\mathcal{P}$ un plan vectoriel euclidien orienté.

    Soient $u, v \in \mathcal{P}\setminus\{0\}$.

    On appelle l'\textbf{angle orienté} de $u$ à $v$, noté $(\widehat{\overrightarrow{u}, \overrightarrow{v}})$,
    le réel $\theta$ défini par le théorème~\ref{theoreme:angle_vecteurs} appliqué à $\frac{u}{\norm{u}}$ et $\frac{v}{\norm{v}}$.
\end{definition}

\begin{remarque}
    L'ensemble des angles orientés de vecteurs du plan est isomorphe à $\R / 2\pi\Z$.
    C'est un groupe abélien pour l'addition.
\end{remarque}

\subsection{Angle de droites}

\begin{proposition}{}{}
    Soient $D$ et $D'$ deux droites d'un plan vectoriel euclidien orienté.

    Si $u_1$ et $u_2$ sont deux vecteurs directeurs de $D$, et $v_1$ et $v_2$ sont deux vecteurs directeurs de $D'$, alors
    $$
    (\widehat{\overrightarrow{u_1}, \overrightarrow{v_1}}) \equiv (\widehat{\overrightarrow{u_2}, \overrightarrow{v_2}}) \mod \pi
    $$
\end{proposition}

\begin{definition}
    Soient $D$ et $D'$ deux droites d'un plan vectoriel euclidien orienté et $u$ et $v$ des vecteurs directeurs de $D$ et $D'$ respectivement.

    On définit l'\textbf{angle orienté} de $D$ à $D'$ par
    $$
    (\widehat{D, D'}) = (\widehat{\overrightarrow{u}, \overrightarrow{v}}) \mod \pi
    $$
\end{definition}

\begin{remarque}
    Par cette définition, on a le resultat peu intuitif de l'égalité $\alpha = \beta$ dans la figure ci-dessous (car $\beta = \alpha + \pi$)
\end{remarque}

\tikzexternalenable
\begin{figure}[H]
    \centering
    \resizebox{0.6\textwidth}{!}{%
    \begin{tikzpicture}[>=stealth, scale=2]

% --- Droites ---
\coordinate (O) at (0,0);
\coordinate (A) at (1.5,0);
\coordinate (B) at (1.5,1.5);
\coordinate (A_neg) at (-1.5,0);
\coordinate (B_neg) at (-1.5,-1.5);

% Droite D
\draw[thick, red] (A_neg) -- (A) node[anchor=east, pos=0.9, fill=white] {$D$};
% Droite D'
\draw[thick, blue] (B_neg) -- (B) node[anchor=south west, pos=0.7, fill=white] {$D'$};

% --- Vecteurs directeurs ---
% Vecteur u sur la demi-droite positive de D
%\draw[->, ultra thick, red!60!black] (O) -- (1,0) node[pos=0.8, yshift=6pt, scale=1.2] {$\mathbf{u}$};
% Vecteur -u sur la demi-droite négative de D
%\draw[->, ultra thick, red!60!black] (O) -- (-1,0) node[pos=0.8, yshift=6pt, scale=1.2] {$-\mathbf{u}$};
% Vecteur v sur la demi-droite positive de D'
%\draw[->, ultra thick, blue!60!black] (O) -- (1,1) node[pos=0.8, xshift=6pt, yshift=3pt, scale=1.2] {$\mathbf{v}$};

% --- Angles ---
% Angle alpha : de u vers v
\pic [draw, ->, angle radius=1.2cm, "$\alpha$", angle eccentricity=1.2, red!80!black] {angle = A--O--B};

% Angle beta : de -u vers v
% Nous traçons l'arc géométrique de la droite D' vers la droite D
\pic [draw, <-, angle radius=1.3cm, "$\beta$", angle eccentricity=1.35, blue!80!black] {angle = B--O--A_neg};

\end{tikzpicture}
    }
\end{figure}
\tikzexternaldisable

\begin{remarque}
    L'ensemble des angles orientés de droites d'un plan est isomorphe à $\R / \pi\Z$.
    C'est un groupe abélien pour l'addition.
\end{remarque}

\begin{remarque}
    On peut caractériser les positions relatives des droites
    \begin{itemize}
        \item $D$ et $D'$ sont parallèles si et seulement si $(\widehat{D, D'}) = 0$.
        \item $D$ et $D'$ sont perpendiculaires si et seulement si $(\widehat{D, D'}) = \frac{\pi}{2}$.
    \end{itemize}
\end{remarque}

\subsection{Chasles et mesures principales}

\begin{theoreme}{Relation de Chasles}{}
    Soit $\mathcal{P}$ un plan vectoriel euclidien orienté.
    \begin{enumerate}
        \item Pour tous $u, v, w \in \mathcal{P}\setminus\{0\}$, on a
        $$
        (\widehat{\overrightarrow{u}, \overrightarrow{v}}) + (\widehat{\overrightarrow{v}, \overrightarrow{w}}) \equiv (\widehat{\overrightarrow{u}, \overrightarrow{w}}) \mod 2\pi
        $$
        \item Pour tous $D, D', D''$ droites d'un plan vectoriel euclidien orienté, on a
        $$
        (\widehat{D, D'}) + (\widehat{D', D''}) \equiv (\widehat{D, D''}) \mod \pi
        $$
    \end{enumerate}
\end{theoreme}

\begin{demonstration}
    Quitte à normaliser les vecteurs, on peut supposer que $u$, $v$ et $w$ sont unitaire.

    Soit $r_1$ la rotation vectorielle telle que $r_1(u) = v$, d'angle $\theta_1 = (\widehat{\overrightarrow{u}, \overrightarrow{v}})$.
    
    Soit $r_2$ la rotation vectorielle telle que $r_2(v) = w$, d'angle $\theta_2 = (\widehat{\overrightarrow{v}, \overrightarrow{w}})$.

    En représentant les rotations vectorielles par des matrices dans la BOND de premier vecteur $u$, on trouve
    $$
    Mat_{\mathcal{B}}(r_1) = \begin{pmatrix}
        \cos(\theta_1) & -\sin(\theta_1) \\
        \sin(\theta_1) & \cos(\theta_1)
    \end{pmatrix}
    $$
    et
    $$
    Mat_{\mathcal{B}}(r_2) = \begin{pmatrix}
        \cos(\theta_2) & -\sin(\theta_2) \\
        \sin(\theta_2) & \cos(\theta_2)
    \end{pmatrix}
    $$
    Ainsi
    $$
    Mat_{\mathcal{B}}(r_2 \circ r_1) = \begin{pmatrix}
        \cos(\theta_1 + \theta_2) & -\sin(\theta_1 + \theta_2) \\
        \sin(\theta_1 + \theta_2) & \cos(\theta_1 + \theta_2)
    \end{pmatrix}
    $$
    Or $r_2 \circ r_1(u) = w$, donc $r_2 \circ r_1$ est la rotation vectorielle d'angle $(\widehat{\overrightarrow{u}, \overrightarrow{w}})$, d'où le résultat sur les vecteurs.

    Le résultat sur les droites s'en déduit en prenant des vecteurs directeurs de ces droites.
\end{demonstration}

\begin{corollaire}{}{}
    Pour tous $u, v \in \mathcal{P}\setminus\{0\}$, on a
    $$
    (\widehat{\overrightarrow{u}, \overrightarrow{v}}) \equiv -(\widehat{\overrightarrow{v}, \overrightarrow{u}}) \mod 2\pi
    $$
\end{corollaire}

\begin{definition}[(Mesure principale d'un angle)]
    Soit $\mathcal{P}$ un plan vectoriel euclidien orienté.

    \begin{itemize}
        \item Pour tous $u, v \in \mathcal{P}\setminus\{0\}$, on appelle \textbf{mesure principale} de l'angle orienté de $u$ à $v$ l'unique réel $\theta$ dans l'intervalle $]-\pi, \pi]$ tel que $$\theta \equiv (\widehat{\overrightarrow{u}, \overrightarrow{v}}) \mod 2\pi$$
        \item Pour tous $D, D'$ droites d'un plan vectoriel euclidien orienté, on appelle \textbf{mesure principale} de l'angle orienté de $D$ à $D'$ l'unique réel $\theta$ dans l'intervalle $]-\frac{\pi}{2}, \frac{\pi}{2}]$ tel que $$\theta \equiv (\widehat{D, D'}) \mod \pi$$
    \end{itemize}
\end{definition}

\subsection{Expression analytique}

\begin{theoreme}{}{}
    Soit $\mathcal{P}$ un plan vectoriel euclidien orienté.

    Soient $u, v \in \mathcal{P}\setminus\{0\}$.

    Alors
    $$
    \langle\overrightarrow{u}, \overrightarrow{v}\rangle = \norm{u}\norm{v}\cos((\widehat{\overrightarrow{u}, \overrightarrow{v}}))
    $$
    et
    $$
    det(\overrightarrow{u}, \overrightarrow{v}) = \norm{u}\norm{v}\sin((\widehat{\overrightarrow{u}, \overrightarrow{v}}))
    $$
\end{theoreme}

\begin{demonstration}
    Posons $i = \dfrac{u}{\norm{u}}$. Il existe un unique $j \in \mathcal{P}$ tel que $\mathcal{B} = (i, j)$ soit une BOND de $\mathcal{P}$.

    Les coordonnées de $u$ dans $\mathcal{B}$ sont $(\norm{u}, 0)$, et $\dfrac{v}{\norm{v}}$ a pour coordonnées $(\cos((\widehat{\overrightarrow{u}, \overrightarrow{v}})), \sin((\widehat{\overrightarrow{u}, \overrightarrow{v}})))$.

    Ainsi $v$ a pour coordonnées $(\norm{v}\cos((\widehat{\overrightarrow{u}, \overrightarrow{v}})), \norm{v}\sin((\widehat{\overrightarrow{u}, \overrightarrow{v}})))$ dans la base $\mathcal{B}$.
    Par conséquent
    $$
    \langle\overrightarrow{u}, \overrightarrow{v}\rangle = \norm{u}\norm{v}\cos((\widehat{\overrightarrow{u}, \overrightarrow{v}}))
    $$
    et
    $$
    det(\overrightarrow{u}, \overrightarrow{v}) = \begin{vmatrix}
        \norm{u} & \norm{v}\cos((\widehat{\overrightarrow{u}, \overrightarrow{v}})) \\
        0 & \norm{v}\sin((\widehat{\overrightarrow{u}, \overrightarrow{v}}))
    \end{vmatrix} = \norm{u}\norm{v}\sin((\widehat{\overrightarrow{u}, \overrightarrow{v}}))
    $$
\end{demonstration}

\subsection{Angle inscrit et cocyclicité}

\begin{lemme}{Angle du triangle isocèle}{}
    Soit $\Omega M A$ un triangle isocèle en $\Omega$, alors
    $$
    2(\widehat{\overrightarrow{MA}, \overrightarrow{M\Omega}}) \equiv (\widehat{\overrightarrow{\Omega A}, \overrightarrow{\Omega M}}) + \pi \mod 2\pi
    $$
\end{lemme}

\begin{demonstration}
    Posons $\theta_1 = (\widehat{\overrightarrow{MA}, \overrightarrow{M\Omega}})$
    et $\theta_2 = (\widehat{\overrightarrow{A \Omega}, \overrightarrow{A M}})$.

    On a
    \begin{align*}
        \langle\overrightarrow{MA}, \overrightarrow{M\Omega}\rangle &= \langle\overrightarrow{M\Omega} + \overrightarrow{\Omega A}, \overrightarrow{M\Omega}\rangle \\
        &= \norm{\overrightarrow{M\Omega}}^2 + \langle\overrightarrow{\Omega A}, \overrightarrow{M\Omega}\rangle
    \end{align*}
    et
    \begin{align*}
        \langle\overrightarrow{A \Omega}, \overrightarrow{A M}\rangle &= \langle-\overrightarrow{\Omega A}, -\overrightarrow{M A}\rangle \\
        &= (-1)(-1)\langle\overrightarrow{\Omega A}, \overrightarrow{M A}\rangle \\
        &= \langle\overrightarrow{\Omega A}, \overrightarrow{M\Omega} + \overrightarrow{\Omega A}\rangle \\
        &= \langle\overrightarrow{\Omega A}, \overrightarrow{M\Omega}\rangle + \norm{\overrightarrow{\Omega A}}^2
    \end{align*}

    Or $\norm{\overrightarrow{M\Omega}} = \norm{\overrightarrow{\Omega A}}$ car le triangle est isocèle en $\Omega$, donc $\cos(\theta_1) = \cos(\theta_2)$

    De même pour les déterminants, on trouve $\sin(\theta_1) = \sin(\theta_2)$.

    Ainsi $(\widehat{\overrightarrow{MA}, \overrightarrow{M\Omega}}) \equiv (\widehat{\overrightarrow{A \Omega}, \overrightarrow{A M}}) \mod 2\pi$.

    La somme des angles d'un triangle est égale à $\pi$ donc
    \begin{align*}  
        &(\widehat{\overrightarrow{M A}, \overrightarrow{M\Omega}}) + (\widehat{\overrightarrow{A \Omega}, \overrightarrow{A M}}) + (\widehat{\overrightarrow{\Omega M}, \overrightarrow{\Omega A}}) \equiv \pi \mod 2\pi \\
        &\implies (\widehat{\overrightarrow{M A}, \overrightarrow{M\Omega}}) + (\widehat{\overrightarrow{A \Omega}, \overrightarrow{A M}}) \equiv (\widehat{\overrightarrow{\Omega A}, \overrightarrow{\Omega M}}) + \pi \mod 2\pi \\
        &\implies 2(\widehat{\overrightarrow{M A}, \overrightarrow{M\Omega}}) \equiv (\widehat{\overrightarrow{\Omega A}, \overrightarrow{\Omega M}}) + \pi \mod 2\pi
    \end{align*}
\end{demonstration}

\pagebreak

\begin{theoreme}{Théorème de l'angle inscrit au centre}{}
    Soit $\Gamma$ un cercle de centre $\Omega$.

    Pour tous $A, B, M \in \Gamma$ distincts, on a
    $$
    2(\widehat{\overrightarrow{MA}, \overrightarrow{MB}}) \equiv (\widehat{\overrightarrow{\Omega A}, \overrightarrow{\Omega B}}) \mod 2\pi
    $$
\end{theoreme}

\tikzexternalenable
\begin{figure}[H]
    \centering
    \resizebox{0.7\textwidth}{!}{%
    \begin{tikzpicture}[>=stealth, scale=1.3]

    % --- Cas 1 : Le centre est à l'intérieur de l'angle ---
    \begin{scope}[shift={(-2.5,0)}]
        \coordinate (O1) at (0,0);
        \draw[thick, gray] (O1) circle (2cm);
        \fill (O1) circle (1.5pt) node[below] {$\Omega$};

        \coordinate (A1) at (220:2cm);
        \coordinate (B1) at (320:2cm);
        \coordinate (M1) at (80:2cm);

        \fill (A1) circle (1.5pt) node[below left] {$A$};
        \fill (B1) circle (1.5pt) node[below right] {$B$};
        \fill (M1) circle (1.5pt) node[above] {$M$};

        \draw[->, thick, red!80!black] (M1) -- (A1);
        \draw[->, thick, red!80!black] (M1) -- (B1);
        \draw[->, thick, blue!80!black] (O1) -- (A1);
        \draw[->, thick, blue!80!black] (O1) -- (B1);

        \pic [draw, ->, "\small $\theta$", angle radius=0.9cm, angle eccentricity=1.3, red!80!black, thick] {angle = A1--M1--B1};
        \pic [draw, ->, "\small $2\theta$", angle radius=0.5cm, angle eccentricity=1.6, blue!80!black, thick] {angle = A1--O1--B1};
        
    \end{scope}

    % --- Cas 2 : Le centre est à l'extérieur de l'angle ---
    \begin{scope}[shift={(2.5,0)}]
        \coordinate (O2) at (0,0);
        \draw[thick, gray] (O2) circle (2cm);
        \fill (O2) circle (1.5pt) node[below left] {$\Omega$};

        % A et B du "même côté" pour que O soit à l'extérieur
        \coordinate (A2) at (10:2cm);
        \coordinate (B2) at (70:2cm);
        \coordinate (M2) at (160:2cm);

        \fill (A2) circle (1.5pt) node[right] {$A$};
        \fill (B2) circle (1.5pt) node[above right] {$B$};
        \fill (M2) circle (1.5pt) node[above left] {$M$};

        \draw[->, thick, red!80!black] (M2) -- (A2);
        \draw[->, thick, red!80!black] (M2) -- (B2);
        \draw[->, thick, blue!80!black] (O2) -- (A2);
        \draw[->, thick, blue!80!black] (O2) -- (B2);

        \pic [draw, ->, "\small $\theta$", angle radius=1.3cm, angle eccentricity=1.2, red!80!black, thick] {angle = A2--M2--B2};
        \pic [draw, ->, "\small $2\theta$", angle radius=0.6cm, angle eccentricity=1.4, blue!80!black, thick] {angle = A2--O2--B2};
    \end{scope}

\end{tikzpicture}
    }
\end{figure}
\tikzexternaldisable

\begin{demonstration}
    On a
    $$
    2(\widehat{\overrightarrow{MA}, \overrightarrow{MB}}) \equiv 2(\widehat{\overrightarrow{MA}, \overrightarrow{M\Omega}}) + 2(\widehat{\overrightarrow{M\Omega}, \overrightarrow{MB}}) \mod 2\pi
    $$

    Or $\Omega M A$ et $\Omega M B$ sont des triangles isocèles en $\Omega$ car $A$, $B$ et $M$ sont sur le cercle de centre $\Omega$.

    Ainsi, en appliquant le lemme précédent
    $$
    2(\widehat{\overrightarrow{MA}, \overrightarrow{M\Omega}}) \equiv (\widehat{\overrightarrow{\Omega A}, \overrightarrow{\Omega M}}) + \pi \mod 2\pi
    $$
    et
    $$
    2(\widehat{\overrightarrow{M\Omega}, \overrightarrow{MB}}) \equiv (\widehat{\overrightarrow{\Omega M}, \overrightarrow{\Omega B}}) + \pi \mod 2\pi
    $$

    Puis par Chasles
    \begin{align*}
        2(\widehat{\overrightarrow{MA}, \overrightarrow{MB}}) &\equiv (\widehat{\overrightarrow{\Omega A}, \overrightarrow{\Omega M}}) + (\widehat{\overrightarrow{\Omega M}, \overrightarrow{\Omega B}}) + 2\pi \mod 2\pi \\
        &\equiv (\widehat{\overrightarrow{\Omega A}, \overrightarrow{\Omega B}}) + 2\pi \mod 2\pi \\
        &\equiv (\widehat{\overrightarrow{\Omega A}, \overrightarrow{\Omega B}}) \mod 2\pi
    \end{align*}
\end{demonstration}

\begin{theoreme}{Condition de cocyclicité}{}
    Soient $A$, $B$, $C$ et $D$ quatre points, trois à trois non alignés, dans un plan vectoriel euclidien orienté.

    Ces quatre points sont cocycliques si et seulement si
    $$
    (\widehat{(CA), (CB)}) \equiv (\widehat{(DA), (DB)}) \mod \pi
    $$
\end{theoreme}

\tikzexternalenable
\begin{figure}[H]
    \centering
    \resizebox{0.9\textwidth}{!}{%
    \begin{tikzpicture}[>=stealth, scale=1.3]

    % --- Cas 1 : Quadrilatère convexe (Carré) ---
    \begin{scope}[shift={(-2.8,0)}]
        \coordinate (O1) at (0,0);
        \draw[dashed, gray] (O1) circle (2cm);

        % Points ordonnés sur le cercle
        \coordinate (A1) at (210:2cm);
        \coordinate (B1) at (330:2cm);
        \coordinate (C1) at (120:2cm);
        \coordinate (D1) at (50:2cm);

        \fill (A1) circle (1.5pt) node[below left] {$A$};
        \fill (B1) circle (1.5pt) node[below right] {$B$};
        \fill (C1) circle (1.5pt) node[above left] {$C$};
        \fill (D1) circle (1.5pt) node[above right] {$D$};

        % Droites prolongées avec calc
        \draw[thick, red!80!black] ($(A1)!-0.2!(C1)$) -- ($(C1)!-0.3!(A1)$);
        \draw[thick, red!80!black] ($(B1)!-0.2!(C1)$) -- ($(C1)!-0.3!(B1)$);

        \draw[thick, blue!80!black] ($(A1)!-0.2!(D1)$) -- ($(D1)!-0.3!(A1)$);
        \draw[thick, blue!80!black] ($(B1)!-0.2!(D1)$) -- ($(D1)!-0.3!(B1)$);

        % Angles
        \pic [draw, ->, angle radius=0.7cm, angle eccentricity=1.4, red!80!black, thick, pic text={\small $\alpha$}] {angle = A1--C1--B1};
        \pic [draw, ->, angle radius=0.7cm, angle eccentricity=1.4, blue!80!black, thick, pic text={\small $\alpha$}] {angle = A1--D1--B1};
        
    \end{scope}

    % --- Cas 2 : Quadrilatère croisé (Papillon) ---
    \begin{scope}[shift={(2.8,0)}]
        \coordinate (O2) at (0,0);
        \draw[dashed, gray] (O2) circle (2cm);

        % D est basculé de l'autre côté de la corde [AB]
        \coordinate (A2) at (210:2cm);
        \coordinate (B2) at (330:2cm);
        \coordinate (C2) at (120:2cm);
        \coordinate (D2) at (250:2cm);

        \fill (A2) circle (1.5pt) node[left] {$A$};
        \fill (B2) circle (1.5pt) node[right] {$B$};
        \fill (C2) circle (1.5pt) node[above left] {$C$};
        \fill (D2) circle (1.5pt) node[below] {$D$};

        % Droites prolongées
        \draw[thick, red!80!black] ($(A2)!-0.2!(C2)$) -- ($(C2)!-0.3!(A2)$);
        \draw[thick, red!80!black] ($(B2)!-0.2!(C2)$) -- ($(C2)!-0.3!(B2)$);

        \draw[thick, blue!80!black] ($(A2)!-0.2!(D2)$) -- ($(D2)!-0.3!(A2)$);
        \draw[thick, blue!80!black] ($(B2)!-0.2!(D2)$) -- ($(D2)!-0.3!(B2)$);

        % Angles
        \pic [draw, ->, angle radius=0.7cm, angle eccentricity=1.4, red!80!black, thick, pic text={\small $\alpha$}] {angle = A2--C2--B2};
        \pic [draw, <-, angle radius=0.4cm, angle eccentricity=1.6, blue!80!black, thick, pic text={\small $\alpha$}] {angle = B2--D2--A2};
        
    \end{scope}

\end{tikzpicture}
    }
\end{figure}
\tikzexternaldisable

\begin{demonstration}
    \begin{description}
        \item[Sens direct.]
        Soit $\Omega$ le centre du cercle passant par $A$, $B$, $C$ et $D$.

        Par le théorème de l'angle inscrit au centre appliqué $A$, $B$, $C$ puis à $A$, $B$, $D$, on trouve
        $$
        2(\widehat{\overrightarrow{CA}, \overrightarrow{CB}}) \equiv (\widehat{\overrightarrow{\Omega A}, \overrightarrow{\Omega B}}) \mod 2\pi
        $$
        et
        $$
        2(\widehat{\overrightarrow{DA}, \overrightarrow{DB}}) \equiv (\widehat{\overrightarrow{\Omega A}, \overrightarrow{\Omega B}}) \mod 2\pi
        $$
        Ainsi
        \begin{align*}
            &2(\widehat{\overrightarrow{CA}, \overrightarrow{CB}}) &\equiv 2(\widehat{\overrightarrow{DA}, \overrightarrow{DB}}) \mod 2\pi \\
            &\iff (\widehat{\overrightarrow{CA}, \overrightarrow{CB}}) &\equiv (\widehat{\overrightarrow{DA}, \overrightarrow{DB}}) \mod \pi \\
            &\iff (\widehat{(CA), (CB)}) &\equiv (\widehat{(DA), (DB)}) \mod \pi
        \end{align*}

        \item[Sens réciproque.]
        Supposons que $(\widehat{(CA), (CB)}) \equiv (\widehat{(DA), (DB)}) \mod \pi$.

        Il existe un unique cercle passant par $A$, $B$ et $C$, de centre $\Omega$.

        Soit $D'$ le point d'intersection de ce cercle avec la droite $(AD)$.

        Puisque $A$, $B$, $C$ et $D'$ sont cocycliques, on a
        \begin{align*}
            &(\widehat{(CA), (CB)}) &\equiv (\widehat{(D'A), (D'B)}) \mod \pi \\
            &\iff (\widehat{(D'A), (D'B)}) &\equiv (\widehat{(DA), (DB)}) \mod \pi \\
            &\iff (\widehat{(D'A), (D'B)}) &\equiv (\widehat{(D'A), (DB)}) \mod \pi \\
            &\iff (\widehat{(D'B), (DB)}) &\equiv 0 \mod \pi \\
        \end{align*}
        Donc $(D'B) = (DB)$ et $D \in (D'A)$ donc $D = D'$, ainsi $D$ est sur le cercle de centre $\Omega$ passant par $A$, $B$ et $C$.
    \end{description}
\end{demonstration}

\end{document}